\section{Conceptos fundamentales}

\textbf{Volumetría (titulación):} Técnica analítica en la que se determina la concentración de una solución (analito) mediante la reacción con otra solución de concentración conocida (titulante).
\textbf{Valoración:} Proceso de añadir el titulante hasta que la reacción se completa (punto de equivalencia).
\textbf{Solución patrón:} Solución de concentración exactamente conocida.
\textbf{Indicador:} Sustancia que cambia de color para indicar el punto final de la titulación.


\section{Tipos de valoraciones}
\begin{itemize}
    \item \textbf{Neutralización (ácido-base):} Reacción entre un ácido y una base.
    \item \textbf{Precipitación:} Formación de un precipitado.
    \item \textbf{Formación de complejos (complejometría):} Formación de un complejo estable (ej. con EDTA).
    \item \textbf{Redox:} Reacción de oxidación-reducción.
\end{itemize}

\section{Valoración complejométrica con EDTA}
El EDTA (ácido etilendiaminotetraacético) es un agente quelante que forma complejos estables con muchos iones metálicos. Se utiliza para determinar la dureza del agua (concentración de $Ca^{2+}$ y $Mg^{2+}$).

\section{Cálculos de concentración}
La normalidad ($N$) es una medida de concentración que relaciona el número de equivalentes de soluto por litro de solución. En las titulaciones, se usa la relación:
\[
N_{1} \times V_{1} = N_{2} \times V_{2}
\]

\resumebox{ejercicio}{Titulación ácido-base}{\footnotesize
    Se titulan 25 mL de una muestra de HCl con NaOH 0.1 N. El punto de equivalencia se alcanza en 20 mL. Calcula la normalidad del HCl.
}

\resumebox{dato}{Dureza del agua}{\footnotesize
    En la industria de alimentos, la dureza del agua es importante porque afecta la textura de los productos, la eficacia de los detergentes y la formación de depósitos en los equipos. La determinación de la dureza se realiza comúnmente por titulación con EDTA.
}